\clearpage\selectlanguage{german}\fancyhead[L]{\small HUNDERT THEOREME ÜBER DIE NEUEN TRANSSCENDENTEN.}\pdfbookmark[0]{Hundert Theoreme über die neuen Transscendenten IV}{transscendentes-iv}\SourceLeaf{470}

\begin{center}
[IV.]\par
\bigskip
\textbf{HUNDERT THEOREME ÜBER DIE NEUEN TRANSSCENDENTEN.}\par
\medskip
\rule{.09\textwidth}{.3pt}\par
\bigskip
\pdfbookmark[1]{Art. 1}{transscendentes-iv-art-1}1.
\end{center}
Es sei
\[
\begin{gathered}
T={}1+\frac{a^n-1}{a-1}.t
+\frac{a^n-1.a^n-a}{a-1.aa-1}.tt
\\[8pt]\quad{}+\frac{a^n-1.a^n-a.a^n-aa}{a-1.aa-1.a^3-1}.t^3\\[8pt]
+\frac{a^n-1.a^n-a.a^n-aa.a^n-a^3}
{a-1.aa-1.a^3-1.a^4-1}.t^4+\text{u.\,s.\,w.}
\end{gathered}
\]
indem wir $n,a,t$ ganz unbestimmt lassen. So oft $n$ eine ganze nicht negative
Zahl ist, bricht die Reihe offenbar ab und besteht aus $n+1$ Gliedern, auch sind
dann, wie wir in der Abhandlung \textit{Summatio quarundam serierum singularium} gezeigt haben, alle Coëfficienten ungebrochene Functionen von $a$. Ist aber $n$ gebrochen oder negativ, so findet beides nicht Statt.

Indem man $T$ mit $1+a^n t$ multiplicirt erhält man
\[
\begin{gathered}T.(1+a^n t)=1+\frac{a^{n+1}-1}{a-1}.t
+\frac{a^{n+1}-1.a^{n+1}-a}{a-1.aa-1}.tt
\\[8pt]\quad{}+\frac{a^{n+1}-1.a^{n+1}-a.a^{n+1}-aa}{a-1.aa-1.a^3-1}.t^3
+\text{etc.}\end{gathered}
\]

Indem man also in $T$ das Element $n$ als veränderlich ansieht, und sich
des Functionalzeichens $\Theta$ bedient, dass
\[
T=\Theta n
\]
wird man haben
\[
\Theta(n+1)=(1+a^n t)\Theta n
\]
Hieraus folgt das 1.\ \textsc{Theorem.}

\SourceLeaf{471}


\begin{center}
\textit{Wenn $n$ eine ganze positive Zahl bedeutet, ist}
\end{center}
\[
\begin{gathered}
1+\frac{a^n-1}{a-1}\,t
 +\frac{a^n-1.a^n-a}{a-1.aa-1}\,tt
 \\[8pt]\quad{}+\frac{a^n-1.a^n-a.a^n-aa}{a-1.aa-1.a^3-1}\,t^3
 +\text{etc.}\ .\ .
\\[8pt]
\qquad=(1+t)(1+at)(1+aat)(1+a^3t)\mathinner{.\,.\,.}(1+a^{n-1}t)
\end{gathered}
\]

\par\medskip\Needspace{5\baselineskip}\pdfbookmark[1]{Art. 2}{transscendentes-iv-art-2}\centerline{2.}\nobreak\smallskip
Wenn wir $T$ auf folgende Art schreiben
\[
\begin{gathered}
1+\frac{1-a^n}{1-a}\,t
 \\[8pt]\quad{}+\frac{1-a^n}{1-a}.\frac{1-a^{n-1}}{1-aa}.att
 \\[8pt]\quad{}+\frac{1-a^n}{1-a}.\frac{1-a^{n-1}}{1-aa}.
  \frac{1-a^{n-2}}{1-a^3}\,a^3t^3
\\[8pt]
\qquad\\[8pt]\quad{}+\frac{1-a^n}{1-a}.\frac{1-a^{n-1}}{1-aa}.
  \frac{1-a^{n-2}}{1-a^3}.\frac{1-a^{n-3}}{1-a^4}.a^6.t^4
 +\text{etc.}
\end{gathered}
\]
wo die Exponenten von $a$ die Trigonalzahlen sein werden, so erhellet,
dass das letzte Glied sein wird
\[
a^{\frac12(nn-n)}t^n=y^n
\]
wenn wir $a^{\frac12(n-1)}t=y$ setzen. Die ganze Reihe wird dann,
indem wir das letzte Glied mit dem ersten, das vorletzte mit dem zweiten
etc. zusammenfassen, für ein gerades $n$
\[
\begin{gathered}
(1+y^n)+\frac{1-a^n}{1-a}\,t(1+y^{n-2})
 \\[8pt]\quad{}+\frac{1-a^n}{1-a}.\frac{1-a^{n-1}}{1-aa}\,at(1+y^{n-4})+.\ .
\\[8pt]
\qquad+\frac{1-a^n}{1-a}.\frac{1-a^{n-1}}{1-aa}..\frac{1-a^{\frac12n+3}}{1-a^{\frac12n-1}}
 a^{\frac12(\frac12n-1)(\frac12n-2)}t^{\frac12n-1}(1+yy)
\\[8pt]
\qquad+\frac{1-a^n}{1-a}.\frac{1-a^{n-1}}{1-aa}..\frac{1-a^{\frac12n+1}}{1-a^{\frac12n}}
 a^{\frac12.\frac12n.\frac12n-1}t^{\frac12n}
\end{gathered}
\]
indem das mittelste Glied isolirt stehen bleibt. Bezeichnen wir dasselbe
durch $A$ und setzen $a=xx$, so wird die Reihe
\[
\begin{gathered}
A\bigl\{1+\frac{1-x^n}{1-x^{n+2}}\,x(y+y^{-1})
 \\[8pt]\quad{}+\frac{1-x^n}{1-x^{n+2}}.\frac{1-x^{n-2}}{1-x^{n+4}}\,x^4(yy+y^{-2})
\\[8pt]
+\frac{1-x^n}{1-x^{n+2}}.\frac{1-x^{n-2}}{1-x^{n+4}}.
 \frac{1-x^{n-4}}{1-x^{n+6}}\,x^9(y^3+y^{-3})+\mathinner{.\,.}\bigr\}
\end{gathered}
\]
u.\,s.\,w. welche Reihe aus $\frac12(n+2)$ Gliedern besteht und dann abbricht.

Unser Product $(1+t)(1+at)(1+aat)..(1+a^{n-1}t)$ hingegen verwandelt sich

